\documentclass[12pt, spanish]{article}
\documentclass{article}
\usepackage[utf8]{inputenc}
\usepackage[T1]{fontenc}
\usepackage[spanish]{babel}
\usepackage{graphicx}
\usepackage{float}
\usepackage{array}
\usepackage{booktabs}
\usepackage{geometry}
\usepackage{fancyhdr}
\usepackage{setspace}
\usepackage{parskip}
\usepackage{amsmath}
\usepackage{amsthm} 
\usepackage{amsfonts}
\usepackage{amssymb}
\usepackage{xcolor}
\usepackage{listings}
\usepackage{caption} 
\usepackage{subcaption}
\usepackage{caption}
\usepackage{hyperref}
\usepackage{geometry}
\geometry{margin=1in}
\usepackage{tikz}
\usepackage{eso-pic}
%\usepackage{siunitx}
\usetikzlibrary{positioning, shapes.geometric, arrows.meta}
\graphicspath{ {./Imagenes/} }
\usepackage[pages=some]{background}

\lstset{
    basicstyle=\ttfamily\scriptsize, % Bajamos de small a scriptsize por lo largo de los #pragma
    breaklines=true,                 % Activa el quiebre de línea
    breakatwhitespace=false,         % Cambia a false para que rompa incluso si no hay espacios
    columns=fullflexible,            % Ayuda a que el espaciado sea más natural
    keepspaces=true,                 % Mantiene los espacios del código original
    literate={á}{{\'a}}1 {é}{{\'e}}1 {í}{{\'i}}1 {ó}{{\'o}}1 {ú}{{\'u}}1 {ñ}{{\~n}}1 % Para que acepte acentos en comentarios
}

%\usepackage[backend=biber,style=authoryear]{biblatex}
%\usepackage{csquotes}
%\addbibresource{referencias.bib}

\definecolor{codegreen}{rgb}{0,0.6,0}
\definecolor{codegray}{rgb}{0.5,0.5,0.5}
\definecolor{codeblue}{rgb}{0,0,0.95}
\definecolor{backcolour}{rgb}{0.95,0.95,0.92}

\lstdefinestyle{mystyle}{
    backgroundcolor=\color{backcolour},   
    commentstyle=\color{codegreen},
    keywordstyle=\color{codeblue},
    numberstyle=\tiny\color{codegray},
    stringstyle=\color{codegreen},
    basicstyle=\ttfamily\footnotesize,
    breakatwhitespace=false,         
    breaklines=true,                 
    captionpos=b,                    
    keepspaces=true,                 
    numbers=left,                    
    numbersep=5pt,                  
    showspaces=false,                
    showstringspaces=false,
    showtabs=false,                  
    tabsize=2
}
\lstset{style=mystyle}

% ============================================================
%  SECCIÓN: Pruebas de Cotas y Argumento de Órbitas
%  Para insertar en reporte.tex después de \subsection{Cotas Fundamentales}
% ============================================================
\begin{document}


\subsection*{Pruebas de las Cotas Fundamentales}
\label{sub:pruebas_cotas}

\subsubsection*{Cota de Hamming (esfera-empaquetamiento)}

\begin{proposition}
Sea $t = \lfloor \frac{d-1}{2} \rfloor$. Entonces
\[
  A(n,d) \;\le\; \frac{2^n}{\displaystyle\sum_{i=0}^{t}\binom{n}{i}}.
\]
\end{proposition}

\begin{proof}
Sea $C \subseteq \{0,1\}^n$ un código con distancia mínima $d$ y $|C| = M$.
Para cada $c \in C$ definimos la \emph{bola cerrada} de radio $t$:
\[
  B(c,\,t) \;=\; \bigl\{\,x \in \{0,1\}^n : d_H(x,c) \le t\,\bigr\}.
\]
Afirmamos que estas bolas son \textbf{disjuntas}.
En efecto, supongamos que existe $x \in B(c,t) \cap B(c',t)$ con $c \neq c'$.
Entonces, por la desigualdad triangular de la distancia de Hamming,
\[
  d_H(c,c') \;\le\; d_H(c,x) + d_H(x,c') \;\le\; t + t \;=\; 2t \;\le\; d-1 < d,
\]
lo cual contradice que $d$ es la distancia mínima de $C$.

Como cada bola contiene exactamente $\sum_{i=0}^{t}\binom{n}{i}$ elementos
(los vectores que difieren de $c$ en exactamente $i$ coordenadas, $0 \le i \le t$),
y todas están contenidas en $\{0,1\}^n$, la condición de disjunción implica
\[
  M \cdot \sum_{i=0}^{t}\binom{n}{i} \;\le\; \bigl|\{0,1\}^n\bigr| = 2^n.
\]
Despejando $M$ obtenemos el resultado.
\end{proof}

\subsubsection*{Cota de Singleton}

\begin{proposition}
$A(n,d) \le 2^{n-d+1}$.
\end{proposition}

\begin{proof}
Sea $C$ un código con distancia mínima $d$.
Considera la proyección $\pi : \{0,1\}^n \to \{0,1\}^{n-d+1}$ que retiene
únicamente las primeras $n-d+1$ coordenadas.
Si $x, y \in C$ con $x \neq y$ satisfacieran $\pi(x) = \pi(y)$,
entonces $x$ e $y$ coincidirían en las primeras $n-d+1$ coordenadas y,
por lo tanto, sólo podrían diferir en las restantes $d-1$ posiciones,
dando $d_H(x,y) \le d-1 < d$, contradicción.
Luego $\pi\!\restriction_C$ es inyectiva, y como
$|\{0,1\}^{n-d+1}| = 2^{n-d+1}$, se sigue $|C| \le 2^{n-d+1}$.
\end{proof}

\subsubsection*{Cota de Plotkin}

\begin{proposition}
Si $2d > n$, entonces $A(n,d) \le \left\lfloor \dfrac{2d}{2d-n} \right\rfloor$.
\end{proposition}

\begin{proof}
Sea $C = \{c_1,\dots,c_M\}$. Contamos la suma de todas las distancias entre pares:
\[
  S = \sum_{1 \le i < j \le M} d_H(c_i, c_j).
\]
\textbf{Cota inferior de $S$.}
Como $d_H(c_i,c_j) \ge d$ para todo par $i \neq j$:
\[
  S \;\ge\; \binom{M}{2} d \;=\; \frac{M(M-1)}{2}\,d.
\]
\textbf{Cota superior de $S$.}
Para cada coordenada $k \in \{1,\dots,n\}$, sea $n_k$ la cantidad de palabras
de $C$ con un $1$ en la posición $k$.
La contribución de la coordenada $k$ a $S$ es $n_k(M - n_k)$.
Usando $n_k(M-n_k) \le M^2/4$:
\[
  S \;=\; \sum_{k=1}^{n} n_k(M-n_k) \;\le\; \frac{nM^2}{4}.
\]
\textbf{Combinando.}
\[
  \frac{M(M-1)}{2}\,d \;\le\; \frac{nM^2}{4}
  \;\implies\; 2d(M-1) \;\le\; nM
  \;\implies\; (2d - n)M \;\le\; 2d.
\]
Como $2d > n$ se tiene $2d - n > 0$, así que $M \le \dfrac{2d}{2d-n}$.
Tomando parte entera: $A(n,d) \le \left\lfloor \dfrac{2d}{2d-n}\right\rfloor$.
\end{proof}

\subsubsection*{Cota de Gilbert–Varshamov (cota inferior)}

\begin{proposition}
$A(n,d) \;\ge\; \dfrac{2^n}{\displaystyle\sum_{i=0}^{d-1}\binom{n}{i}}$.
\end{proposition}

\begin{proof}
Construimos $C$ de manera \emph{greedy}: partiendo de $C = \emptyset$,
en cada paso agregamos a $C$ una palabra de $\{0,1\}^n$ que tenga distancia
$\ge d$ con todas las palabras ya incluidas.
El proceso termina cuando ninguna palabra restante puede ser añadida,
es decir, cuando cada $x \in \{0,1\}^n \setminus C$ cumple $d_H(x,c) < d$
para algún $c \in C$; en otras palabras, $x \in B(c, d-1)$ para algún $c \in C$.
Esto significa
\[
  \{0,1\}^n \;\subseteq\; \bigcup_{c \in C} B(c,\,d-1),
\]
y tomando cardinalidades:
\[
  2^n \;\le\; |C| \cdot \sum_{i=0}^{d-1}\binom{n}{i}.
\]
Despejando obtenemos $|C| \ge \dfrac{2^n}{\sum_{i=0}^{d-1}\binom{n}{i}}$,
y como $C$ es un código válido, $A(n,d)$ es al menos tan grande.
\end{proof}

\end{document}